\section{Related Work}
\label{sec:related}

\textbf{Response-type measurement and physical vehicle models.}
IRI is defined on the longitudinal profile through a reference quarter car (the ``Golden Car'') driven at 80\,km/h~\cite{sayers1986tp45,sayers1995iri,astm_e1926}.
Response-type systems measure a vehicle's reaction instead of the profile and must be correlated with a reference profiler for each vehicle~\cite{sayers1986tp46}.
Model-based methods estimate roughness from the responses of a specific vehicle with fixed or identified vehicle models and transfer functions~\cite{gonzalez2008accel,zhao2017freq}; Bridgelall's road impact factor normalises vertical-acceleration energy for connected vehicles~\cite{bridgelall2014rif}.
Wu et al.~\cite{wu2023calib} update a full-car model of the test vehicle from a few impact tests, derive a linear relation between IRI and vibration RMS by simulation over random profiles, and validate it on real roads; this calibrates the vehicle model instead of the estimator.
Closest to our data, Skar and Levenberg~\cite{skar2024inversion} invert road profiles from the accelerations of ten nominally identical electric cars with a quarter-car model, a PID controller and nonlinear error minimisation, and compare them with a laser profilometer; Agebj\"ar et al.~\cite{agebjar2025fusion} identify a vehicle model of a production car and reconstruct the profile with a Kalman filter.
These methods need a model of each vehicle with known or identified parameters; our baseline M7 follows the same idea with a quarter car calibrated on the training routes, but it is not a reproduction of~\cite{skar2024inversion}, whose calibrated vehicle parameters are not available (its open code contains only the Golden-Car parameters).

\textbf{Learned estimators on real data.}
Smartphone studies report agreement with reference roughness that depends on speed, vehicle and phone~\cite{thilak2021}, on mounting~\cite{mounting2026} and on the number of repeated runs~\cite{opara2022roughness}; see also~\cite{hanson2014iri,sattar2018review,evalframework2023}.
Data-driven models accounting for speed, vehicle type and mounting were trained on real smartphone data~\cite{sang2024iri}, semi-supervised learning exploits unlabelled responses of crowdsourced vehicles~\cite{zhang2026semisup,trc2021crowd}, and Jeong and Jo~\cite{jeongjo2024irinet} validated a deep IRI estimator on data of 29 vehicles against an inertial profiler.
On LiRA data, Bajic et al.~\cite{bajic2021ml} trained regressors on the car's vertical acceleration and speed along the M3 motorway and tested them on the last 20\,\% of the same road; we test on routes not used for training and report uncertainty that accounts for spatial dependence.

\textbf{Learning from simulated responses.}
Closest in design to our study, Nitsche et al.~\cite{nitsche2014wlp} trained regressors on features of simulated vehicle responses and validated them with a probe vehicle, for the weighted longitudinal profile rather than IRI (according to Bajic et al.~\cite{bajic2021ml}, over a measured profile, against an inertial profilometer and without varying speed).
Networks trained on simulated responses have also been used to reconstruct profiles and classify roughness~\cite{ngwangwa2010ann}, then, trained on a full-vehicle model over slightly modified measured profiles, applied to that vehicle's measured responses~\cite{ngwangwa2014exp}, to classify roughness grades from quarter-car simulations~\cite{rcnn2023sim} and, with a real-vehicle test, from the unsprung-mass response~\cite{speedadaptive2025}, and to estimate IRI without vehicle-specific calibration: Jeong et al.~\cite{jeong2020cnn} trained CNNs on simulated responses of four vehicle types with real driving-speed profiles and validated them numerically with real road IRI information.
Randomising the simulator~\cite{tobin2017dr,peng2018dr} and adapting the randomisation with a few real roll-outs~\cite{chebotar2019simopt} are the standard devices for simulation-to-real transfer; calibration or joint training with paired drives and the same-segment spectral comparison play this role here.
Game engines provide vehicle dynamics for software testing~\cite{gambi2019beamng,beamng}, but IRI labels derived from their authored road meshes need validation (Supplement~\ref{S-sec:s_earlier}).
We add, for continuous IRI from a production car on held-out roads of public data, a decomposition of what a simulation-trained estimator transfers, which simulator choices set its level, and how much paired data must accompany the simulation, for several ways of combining simulated and real data.
